\chapter{Постановка и цел на работата}

Тук обобщавам резултатите от \cite{LinDim}\\
Нека разгледаме диференциалното включване в $\R^n$
\[
    (1) \left|
    \begin{array}{ll}
        \dot{x}(t)\in F(x(t),t) \quad \text{почти навсякъде}\\
        x(0) = x_0 \\
        t \in [0, T],
    \end{array}
    \right.
\]
където $F:\R^n \times \R \rightrightarrows \R^n$ е многозначно изображение непрекъснато по $x$ и измеримо по $t$.\\
Решенията на това уравнение са абсолютно непрекъснати функции $x\in AC([0,T], \R^n) \subseteq L^{\infty}([0,T], \R^n)=X$ с производни $\dot{x}\in L^1([0,T], \R^n)=Y$ изпълняващи горните условия.\\
Да запишем $Q=\{(x,y)\in\Univ \mid x \text{ е решение на (1) и } y \text{ е неговата производна}\}$.\\
Нашата цел ще бъде да изведем някаква информация за допирателния конус на Кларк към множеството $Q$ в точка $\Point\in Q$ - $\Tangent_Q(\Point)$.\\

Ще ползваме представянето $Q = A \cap B$, където:
\[
    A = \left\{(x,y)\in\Univ \mid x(t)=x_0 + \int_0^t y(s)ds\right\}
\]
\[
    B = \{(x,y)\in\Univ \mid y(t)\in F(x(t),t)\ \text{п.н.}\}
\]
или еквивалентно
\[
    B = \{(x,y)\in\Univ \mid (x(t),y(t))\in\Gr\ \text{п.н.}\}.
\]
За удобство ще вземем $x_0 = 0$, но общият случай се обобщава тривиално от нашия.\\

Крайният ни резултат е следният:\\
За многозначни изображения липшицови по $x$ и измерими по $t$:
\[
    \Tangent_Q\Point \supseteq A \cap \{\Vec\in\Univ \mid \Vect\in T_{\Gr}(\Pointt)\ \text{п.н.}\}
\]
